\chapter{Mimari Tasarım}

\section{Genel Mimari Yaklaşım}
Sistem, mikroservis mimarisine dayalı olarak tasarlanmıştır. Bu mimari, sistemin ölçeklenebilirliğini, dayanıklılığını ve bakım kolaylığını artırmak için tercih edilmiştir. Her mikroservis, bağımsız olarak geliştirilebilir, test edilebilir, dağıtılabilir ve ölçeklendirilebilir.

\subsection{Mimari Prensipler}
\begin{itemize}
    \item \textbf{Bağımsız Dağıtılabilirlik:} Her servis bağımsız olarak dağıtılabilir.
    \item \textbf{Tek Sorumluluk:} Her servis belirli bir iş alanından (domain) sorumludur.
    \item \textbf{Veritabanı Ayrımı:} Her servis kendi veritabanına sahiptir (Database per Service).
    \item \textbf{API Tabanlı İletişim:} Servisler REST API ve message queue ile iletişim kurar.
    \item \textbf{Hata Toleransı:} Bir servisin çökmesi tüm sistemi etkilemez.
\end{itemize}

\section{Mikroservisler ve Sorumlulukları}

\subsection{User Service}
\textbf{Sorumluluklar:}
\begin{itemize}
    \item Kullanıcı kaydı ve doğrulama
    \item Kullanıcı girişi ve kimlik doğrulama (JWT token oluşturma)
    \item Profil yönetimi (görüntüleme, güncelleme)
    \item Kullanıcı bilgilerini diğer servislere sağlama
\end{itemize}

\textbf{Teknolojiler:}
\begin{itemize}
    \item Veritabanı: PostgreSQL (Kullanıcı bilgileri)
    \item Bağımlılıklar: Redis (Cache)
\end{itemize}

\subsection{Chat Service}
\textbf{Sorumluluklar:}
\begin{itemize}
    \item Birebir ve grup sohbeti oluşturma
    \item Sohbet listesi yönetimi
    \item Sohbet katılımcı yönetimi (ekleme, çıkarma)
    \item Sohbet metadata yönetimi
\end{itemize}

\textbf{Teknolojiler:}
\begin{itemize}
    \item Veritabanı: MongoDB (Sohbet bilgileri ve katılımcıları)
    \item Bağımlılıklar: User Service, RabbitMQ, Redis
\end{itemize}

\subsection{Message Service}
\textbf{Sorumluluklar:}
\begin{itemize}
    \item Mesaj gönderme ve depolama
    \item Mesaj geçmişi sorgulama (pagination)
    \item Mesaj silme (soft delete)
    \item Mesaj durumu yönetimi (sent, delivered, read, deleted)
\end{itemize}

\textbf{Teknolojiler:}
\begin{itemize}
    \item Veritabanı: MongoDB (Mesaj bilgileri)
    \item Bağımlılıklar: RabbitMQ, Redis, User Service
\end{itemize}

\subsection{Notification Service}
\textbf{Sorumluluklar:}
\begin{itemize}
    \item Gerçek zamanlı bildirim gönderme (WebSocket)
    \item Bildirim kuyruğu yönetimi (çevrimdışı kullanıcılar için)
    \item Bildirim geçmişi saklama
    \item Bildirim durumu yönetimi (okundu/okunmadı)
\end{itemize}

\textbf{Teknolojiler:}
\begin{itemize}
    \item Veritabanı: MongoDB (Bildirim bilgileri)
    \item Bağımlılıklar: RabbitMQ, WebSocket Gateway, Redis, Chat Service
\end{itemize}

\subsection{File Service}
\textbf{Sorumluluklar:}
\begin{itemize}
    \item Dosya yükleme ve depolama
    \item Dosya indirme ve URL oluşturma
    \item Dosya türü ve boyut kontrolü
    \item Güvenli dosya erişimi
\end{itemize}

\textbf{Teknolojiler:}
\begin{itemize}
    \item Veritabanı: PostgreSQL (Dosya metadata)
    \item Bağımlılıklar: Storage Service (Local file system)
\end{itemize}

\section{Altyapı Bileşenleri}

\subsection{API Gateway}
Tüm dış HTTP isteklerini yönlendirir, kimlik doğrulama, yetkilendirme ve rate limiting (60 request/dakika/IP) işlemlerini gerçekleştirir. Node.js (Express) ve http-proxy-middleware kullanılmıştır.

\subsection{WebSocket Gateway}
Gerçek zamanlı mesaj iletimi ve bağlantı yönetimini sağlar. Kullanıcı bağlantı durumu Redis üzerinde takip edilir. Node.js ve Socket.io kullanılmıştır.

\subsection{Message Queue (RabbitMQ)}
Servisler arası asenkron iletişimi ve event-driven mimariyi destekler. `chat.exchange` ve `message.exchange` üzerinden bildirimler Notification Service'e iletilir.

\subsection{Cache (Redis)}
Sık kullanılan verilerin (sohbet listesi, son mesajlar) cache'lenmesi ve kullanıcı bağlantı durumunun takibi için kullanılır.

\subsection{Veritabanları}
Database per Service stratejisi uygulanmıştır:
\begin{itemize}
    \item \textbf{PostgreSQL:} User Service ve File Service (İlişkisel veriler)
    \item \textbf{MongoDB:} Chat Service, Message Service ve Notification Service (Doküman tabanlı, yüksek yazma hızı gerektiren veriler)
\end{itemize}

\section{Servisler Arası İletişim Desenleri}

\subsection{Senkron İletişim (REST API)}
Anında yanıt gerektiren işlemler için kullanılır (Örn: Chat Service $\rightarrow$ User Service). Güvenlik için JWT ve servis token'ları kullanılır.

\subsection{Asenkron İletişim (Message Queue)}
Event-driven işlemler ve loose coupling için kullanılır (Örn: Message Service $\rightarrow$ RabbitMQ $\rightarrow$ Notification Service).

\subsection{Gerçek Zamanlı İletişim (WebSocket)}
Anında bildirim ve çift yönlü iletişim için kullanılır (Örn: Notification Service $\rightarrow$ WebSocket Gateway $\rightarrow$ Client).

\section{Teknoloji Stack}

\begin{itemize}
    \item \textbf{Backend:} Node.js, Express.js
    \item \textbf{Frontend:} React 19, TypeScript, Vite, Tailwind CSS, Radix UI
    \item \textbf{Veritabanı:} PostgreSQL 15, MongoDB 7
    \item \textbf{Cache \& Queue:} Redis 7, RabbitMQ 3
    \item \textbf{Altyapı:} Docker, Docker Compose, ELK Stack (Logging)
\end{itemize}

\section{Güvenlik Mimarisi}
\begin{itemize}
    \item \textbf{Kimlik Doğrulama:} JWT Token
    \item \textbf{Veri Güvenliği:} HTTPS, Bcrypt (Şifre hashleme)
    \item \textbf{Kontroller:} Rate Limiting, SQL Injection Koruması, XSS Koruması
\end{itemize}

\section{Ölçeklenebilirlik Stratejisi}
\begin{itemize}
    \item \textbf{Yatay Ölçeklenebilirlik:} Stateless servisler, Load Balancer
    \item \textbf{Veritabanı:} Read Replicas, Sharding
    \item \textbf{WebSocket:} Redis adapter ile çoklu instance desteği
\end{itemize}
